\section{Related Work}
\label{sec:related}

\paragraph{Memory poisoning attacks.}
\citet{chen2024agentpoison} introduce \agentpoison{}, the first attack to
exploit the retrieval mechanism of LLM agent memory systems.  They formulate
trigger optimization as a constrained gradient problem over the embedding
space of a dense retriever (DPR~\citep{karpukhin2020dense}) and demonstrate
attacks on three production agent frameworks.  \citet{dong2025minja} extend
the threat model to a \emph{query-only} setting: the attacker has no direct
memory access and relies on the agent's auto-storage of adversarial reasoning
chains.  Their progressive-shortening indication protocol achieves ISR of
98.2\% across three experimental platforms.  \citet{injecmem2026} propose a
retriever-agnostic anchor strategy that achieves broad recall via
topic-covering passages designed to match many query types.

\paragraph{Backdoor attacks on neural retrieval.}
Retrieval-augmented generation (RAG) systems have been shown to be vulnerable
to corpus poisoning \citep{zou2024poisonedrag}, where adversarial documents
are injected into the retrieval corpus.  Our work differs in that the
adversarial entries are \emph{agent-generated memories} rather than
external documents, and the threat model spans multiple sessions.

\paragraph{Watermark-based content provenance.}
\citet{zhao2024watermark} propose a provably robust token-level watermark
for LLM-generated text based on a green/red vocabulary partition.  We adapt
this scheme to character-level memory entries (Section~\ref{sec:defenses}).
Related work on watermark evasion \citep{kirchenbauer2023watermark} motivates
our evaluation of paraphrasing and copy-paste dilution attacks.

\paragraph{Backdoor defense strategies.}
Neural Cleanse \citep{wang2019neuralcleanse} and STRIP \citep{gao2019strip}
are representative backdoor detection methods that inspired our
attack-defense matrix evaluation protocol.  Unlike these methods, which
operate on model weights, our defenses target memory content directly.
The ASB benchmark \citep{zhang2024asb} provides a unified evaluation
framework for backdoor attacks and defenses in classification settings;
we adapt its matrix format to the memory-agent threat model.

\paragraph{LLM agent security.}
Prompt injection attacks \citep{perez2022ignore} target LLM agents through
crafted user inputs.  Memory poisoning is orthogonal: the adversarial content
is stored silently and triggers on future queries without further adversarial
interaction.
